% --- CAPÍTULO 2: REFERENCIAL TEÓRICO E EPISTEMOLÓGICO ---
\chapter{Enquadramento Teórico-Metodológico}
\label{cap:referencial}

\section{O Regime Discursivo da Ciência}
A construção da autoridade científica repousa historicamente sobre imperativos de impessoalidade, univocidade conceitual e rastreabilidade metodológica...

\section{O Modelo Totem: Máscara, Marca e Mancha}
Para compreender os matizes dessa fricção discursiva, propõe-se um enquadramento analítico triádico fundamentado na estrutura do \textit{Totem}, articulado a partir de três dimensões interdependentes:

\begin{itemize}
    \item \textbf{A Máscara:} Representa a persona institucionalizada do especialista, a performance de neutralidade e o papel social que legitima o sujeito de enunciação do saber.
    \item \textbf{A Marca:} Configura o \textit{pólo articulado} do sistema. É a inscrição rigorosa, o registro textual estável, a codificação metodológica que carimba o discurso com o selo da cientificidade oficial.
    \item \textbf{A Mancha:} Configura o \textit{pólo não articulado}. Compreende o ruído, a opacidade, a complexidade inerente à oralidade, à ancestralidade e à vivência material que resiste à simplificação e à domesticação pelas categorias rígidas da academia.
\end{itemize}